% GENERATED FILE. Edit canonical lessons, then run npm run generate:books.
% canonical-source-hash: fnv1a64:d40e669baeb54453
% canonical-lessons: HI-C405-gati, HI-C405-gaddha, HI-C405-kandaktar, HI-C405-uchalnaintr, HI-C405-kattna, HI-R405-review-398-405

\chapter{Speed and Bumps}
\label{ch:hi-speed-and-bumps}
\hlchaptermodality{405}

\begin{chapteropening}
\textbf{By the end of this chapter:} \emph{I can talk about speed, potholes and the conductor, a bumpy ride and a road that branches off.}

Complete the last lesson of chapter 405: I can talk about speed, potholes and the conductor, a bumpy ride and a road that branches off.
\end{chapteropening}

\section[\texorpdfstring{\hi{गति} (gati)}{गति (gati)}]{\hi{गति} (gati) — speed (formal); motion}

\label{lesson:HI-C405-gati}



\begin{quote}
Before the new one: say the Hindi for to crawl, then the Hindi for to drive (animals).
\end{quote}

\subsection*{You'll want to know: \hi{गति}}
\textbf{\hi{गति}} — \emph{gati} — \textquotedblleft{}speed (formal); motion\textquotedblright{}.

\begin{quote}
\hi{इस} \hi{सड़क} \hi{पर} \hi{गति} \hi{सीमा} \hi{चालीस} \hi{है।}
\end{quote}

\begin{quote}
\emph{is saṛak par gati sīmā cālīs hai.}
\end{quote}

\begin{quote}
The speed limit on this road is forty.
\end{quote}

Feminine.

\begin{cousinweb}
\textbf{\hi{गति}} (\emph{gati}) is a tatsama, Sanskrit \emph{gati}, \textquotedblleft{}going, moving, motion\textquotedblright{} (Monier-Williams), kept unchanged; road signs and the news use it: \hi{गति} \hi{सीमा}, the speed limit.
\end{cousinweb}

\begin{grammarlens}[title={What you've built}]
One more word for getting around.
\end{grammarlens}

\subsection*{Your turn}
\begin{itemize}
\raggedright
  \item \emph{Say it:} \emph{gati}
  \item \emph{Say it:} \emph{gati}, once more
  \item \emph{Say it:} \emph{is saṛak par gati sīmā cālīs hai.} — The speed limit on this road is forty.
  \item \emph{Recall:} say the Hindi for to crawl, then the Hindi for to drive (animals), then say \emph{gati} again
\end{itemize}

\begin{tcolorbox}[breakable,colback=teal!4,colframe=teal!35!black,title={Before you move on}]
What does \hi{गति} mean? (\textquotedblleft{}Speed (formal); motion\textquotedblright{}.) Say it once more.
\end{tcolorbox}
\section[\texorpdfstring{\hi{गड्ढा} (gaḍḍhā)}{गड्ढा (gaḍḍhā)}]{\hi{गड्ढा} (gaḍḍhā) — a pothole; a pit}

\label{lesson:HI-C405-gaddha}



\begin{quote}
Before the new one: say the Hindi for to drive (animals), then the Hindi for speed (formal).
\end{quote}

\subsection*{You'll want to know: \hi{गड्ढा}}
\textbf{\hi{गड्ढा}} — \emph{gaḍḍhā} — \textquotedblleft{}a pothole; a pit\textquotedblright{}.

\begin{quote}
\hi{इस} \hi{सड़क} \hi{में} \hi{बहुत} \hi{गड्ढे} \hi{हैं।}
\end{quote}

\begin{quote}
\emph{is saṛak meṁ bahut gaḍḍhe haiṁ.}
\end{quote}

\begin{quote}
This road has a lot of potholes.
\end{quote}

Masculine.

\begin{cousinweb}
\textbf{\hi{गड्ढा}} (\emph{gaḍḍhā}) belongs with \emph{*gaḍḍa}, \textquotedblleft{}hole, pit\textquotedblright{}, a form CDIAL reconstructs (CDIAL 3981); CDIAL gives Hindi \emph{gāṛā}, \textquotedblleft{}hole\textquotedblright{}, there, but not this form with ḍḍh, so the link is probable rather than shown.
\end{cousinweb}

\begin{grammarlens}[title={What you've built}]
One more word for getting around.
\end{grammarlens}

\subsection*{Your turn}
\begin{itemize}
\raggedright
  \item \emph{Say it:} \emph{gaḍḍhā}
  \item \emph{Say it:} \emph{gaḍḍhā}, once more
  \item \emph{Say it:} \emph{is saṛak meṁ bahut gaḍḍhe haiṁ.} — This road has a lot of potholes.
  \item \emph{Recall:} say the Hindi for to drive (animals), then the Hindi for speed (formal), then say \emph{gaḍḍhā} again
\end{itemize}

\begin{tcolorbox}[breakable,colback=teal!4,colframe=teal!35!black,title={Before you move on}]
What does \hi{गड्ढा} mean? (\textquotedblleft{}A pothole; a pit\textquotedblright{}.) Say it once more.
\end{tcolorbox}
\section[\texorpdfstring{\hi{कंडक्टर} (kaṇḍakṭar)}{कंडक्टर (kaṇḍakṭar)}]{\hi{कंडक्टर} (kaṇḍakṭar) — a conductor; a bus conductor}

\label{lesson:HI-C405-kandaktar}



\begin{quote}
Before the new one: say the Hindi for speed (formal), then the Hindi for a pothole.
\end{quote}

\subsection*{You'll want to know: \hi{कंडक्टर}}
\textbf{\hi{कंडक्टर}} — \emph{kaṇḍakṭar} — \textquotedblleft{}a conductor; a bus conductor\textquotedblright{}.

\begin{quote}
\hi{कंडक्टर} \hi{ने} \hi{किराया} \hi{माँगा।}
\end{quote}

\begin{quote}
\emph{kaṇḍakṭar ne kirāyā mā̃gā.}
\end{quote}

\begin{quote}
The conductor asked for the fare.
\end{quote}

Masculine or feminine, by the person meant.

\begin{cousinweb}
\textbf{\hi{कंडक्टर}} (\emph{kaṇḍakṭar}) is borrowed from English \emph{conductor}, the person who sells tickets on a bus.
\end{cousinweb}

\begin{grammarlens}[title={What you've built}]
One more word for getting around.
\end{grammarlens}

\subsection*{Your turn}
\begin{itemize}
\raggedright
  \item \emph{Say it:} \emph{kaṇḍakṭar}
  \item \emph{Say it:} \emph{kaṇḍakṭar}, once more
  \item \emph{Say it:} \emph{kaṇḍakṭar ne kirāyā mā̃gā.} — The conductor asked for the fare.
  \item \emph{Recall:} say the Hindi for speed (formal), then the Hindi for a pothole, then say \emph{kaṇḍakṭar} again
\end{itemize}

\begin{tcolorbox}[breakable,colback=teal!4,colframe=teal!35!black,title={Before you move on}]
What does \hi{कंडक्टर} mean? (\textquotedblleft{}A conductor; a bus conductor\textquotedblright{}.) Say it once more.
\end{tcolorbox}
\section[\texorpdfstring{\hi{उछलना} (uchalnā)}{उछलना (uchalnā)}]{\hi{उछलना} (uchalnā) — to bounce; to jump up}

\label{lesson:HI-C405-uchalnaintr}



\begin{quote}
Before the new one: say the Hindi for a pothole, then the Hindi for a conductor.
\end{quote}

\subsection*{You'll want to know: \hi{उछलना}}
\textbf{\hi{उछलना}} — \emph{uchalnā} — \textquotedblleft{}to bounce; to jump up\textquotedblright{}.

\begin{quote}
\hi{गड्ढे} \hi{में} \hi{बस} \hi{उछली} \hi{और} \hi{सब} \hi{लोग} \hi{हँस} \hi{पड़े।}
\end{quote}

\begin{quote}
\emph{gaḍḍhe meṁ bas uchlī aur sab log hãs paṛe.}
\end{quote}

\begin{quote}
The bus bounced in the pothole and everyone burst out laughing.
\end{quote}

Past (he …): \hi{वह} \hi{बच्चा} \hi{ख़ुशी} \hi{से} \hi{उछला}, \emph{vah baccā khuśī se uchlā}.

\begin{cousinweb}
\textbf{\hi{उछलना}} (\emph{uchalnā}) goes back to \emph{*ut-śalati}, \textquotedblleft{}springs up\textquotedblright{}, a form CDIAL reconstructs (CDIAL 1843), which lists Hindi \emph{uchalnā}, \textquotedblleft{}to leap\textquotedblright{}.
\end{cousinweb}

\begin{grammarlens}[title={What you've built}]
One more word for getting around.
\end{grammarlens}

\subsection*{Your turn}
\begin{itemize}
\raggedright
  \item \emph{Say it:} \emph{uchalnā}
  \item \emph{Say it:} \emph{uchalnā}, once more
  \item \emph{Say it:} \emph{gaḍḍhe meṁ bas uchlī aur sab log hãs paṛe.} — The bus bounced in the pothole and everyone burst out laughing.
  \item \emph{Recall:} say the Hindi for a pothole, then the Hindi for a conductor, then say \emph{uchalnā} again
\end{itemize}

\begin{tcolorbox}[breakable,colback=teal!4,colframe=teal!35!black,title={Before you move on}]
What does \hi{उछलना} mean? (\textquotedblleft{}To bounce; to jump up\textquotedblright{}.) Say it once more.
\end{tcolorbox}
\section[\texorpdfstring{\hi{कटना} (kaṭnā)}{कटना (kaṭnā)}]{\hi{कटना} (kaṭnā) — to be cut; to branch off (of a road)}

\label{lesson:HI-C405-kattna}



\begin{quote}
Before the new one: say the Hindi for a conductor, then the Hindi for to bounce.
\end{quote}

\subsection*{You'll want to know: \hi{कटना}}
\textbf{\hi{कटना}} — \emph{kaṭnā} — \textquotedblleft{}to be cut; to branch off (of a road)\textquotedblright{}.

\begin{quote}
\hi{आगे} \hi{से} \hi{एक} \hi{सड़क} \hi{बाएँ} \hi{कटती} \hi{है।}
\end{quote}

\begin{quote}
\emph{āge se ek saṛak bāẽ kaṭtī hai.}
\end{quote}

\begin{quote}
Further on, a road branches off to the left.
\end{quote}

Past (it …): \hi{वह} \hi{पतंग} \hi{कट} \hi{गई}, \emph{vah patang kaṭ gaī}.

\hi{बिजली} \hi{कट} \hi{गई}, the power went off.

\begin{cousinweb}
\textbf{\hi{कटना}} (\emph{kaṭnā}) is the intransitive partner of \hi{काटना}, both from Sanskrit \emph{kartati}, \textquotedblleft{}cuts\textquotedblright{} (CDIAL 2854), which lists Hindi \emph{kaṭnā}, \textquotedblleft{}to be cut\textquotedblright{}. A side road \hi{कटती} where it cuts away from the main one.
\end{cousinweb}

\begin{grammarlens}[title={What you've built}]
One more word for getting around.
\end{grammarlens}

\subsection*{Your turn}
\begin{itemize}
\raggedright
  \item \emph{Say it:} \emph{kaṭnā}
  \item \emph{Say it:} \emph{kaṭnā}, once more
  \item \emph{Say it:} \emph{āge se ek saṛak bāẽ kaṭtī hai.} — Further on, a road branches off to the left.
  \item \emph{Recall:} say the Hindi for a conductor, then the Hindi for to bounce, then say \emph{kaṭnā} again
\end{itemize}

\begin{tcolorbox}[breakable,colback=teal!4,colframe=teal!35!black,title={Before you move on}]
What does \hi{कटना} mean? (\textquotedblleft{}To be cut; to branch off (of a road)\textquotedblright{}.) Say it once more.
\end{tcolorbox}
\section[(dialogue)]{Looking back: from Passengers and Porters to Speed and Bumps}

\label{lesson:HI-R405-review-398-405}



\begin{quote}
Say the Hindi for a traveller and a porter.
\end{quote}

\subsection*{Your turn}
Say these aloud:

\begin{itemize}
\raggedright
  \item a traveller, a porter, to haul, to load (onto), to seat (someone)
  \item a line, an engine, a sleeper carriage, a whistle, to slide
  \item a reservation (of a seat), a berth, waiting, to check, setting off
  \item a gate, a signal, a junction, to step over, to look out (at)
  \item a bundle, a suitcase, to slide (something), to doze, to rest a while
  \item a stand, a passer-by, a pavement, to take a stroll, to pant
  \item traffic, a traffic jam, a motorcycle, to crawl, to drive (animals)
  \item speed (formal), a pothole, a conductor, to bounce, to be cut
\end{itemize}

\begin{tcolorbox}[breakable,colback=teal!4,colframe=teal!35!black,title={Before you move on}]
A traveller? (\textbf{\hi{मुसाफ़िर}}, \emph{musāfir}.) A porter? (\textbf{\hi{कुली}}, \emph{kulī}.) To haul? (\textbf{\hi{ढोना}}, \emph{ḍhonā}.) To load (onto)? (\textbf{\hi{लादना}}, \emph{lādnā}.) To seat (someone)? (\textbf{\hi{बिठाना}}, \emph{biṭhānā}.) A line? (\textbf{\hi{लाइन}}, \emph{lāin}.) An engine? (\textbf{\hi{इंजन}}, \emph{injan}.) A sleeper carriage? (\textbf{\hi{स्लीपर}}, \emph{slīpar}.) A whistle? (\textbf{\hi{सीटी}}, \emph{sīṭī}.) To slide? (\textbf{\hi{सरकना}}, \emph{saraknā}.) A reservation (of a seat)? (\textbf{\hi{आरक्षण}}, \emph{ārakṣaṇ}.) A berth? (\textbf{\hi{बर्थ}}, \emph{barth}.) Waiting? (\textbf{\hi{प्रतीक्षा}}, \emph{pratīkṣā}.) To check? (\textbf{\hi{जाँचना}}, \emph{jā̃cnā}.) Setting off? (\textbf{\hi{रवाना}}, \emph{ravānā}.) A gate? (\textbf{\hi{फाटक}}, \emph{phāṭak}.) A signal? (\textbf{\hi{सिग्नल}}, \emph{signal}.) A junction? (\textbf{\hi{जंक्शन}}, \emph{jaṅkśan}.) To step over? (\textbf{\hi{लाँघना}}, \emph{lā̃ghnā}.) To look out (at)? (\textbf{\hi{ताकना}}, \emph{tāknā}.) A bundle? (\textbf{\hi{गठरी}}, \emph{gaṭhrī}.) A suitcase? (\textbf{\hi{अटैची}}, \emph{aṭaicī}.) To slide (something)? (\textbf{\hi{सरकाना}}, \emph{sarkānā}.) To doze? (\textbf{\hi{ऊँघना}}, \emph{ū̃ghnā}.) To rest a while? (\textbf{\hi{सुस्ताना}}, \emph{sustānā}.) A stand? (\textbf{\hi{स्टैंड}}, \emph{sṭaiṇḍ}.) A passer-by? (\textbf{\hi{राहगीर}}, \emph{rāhgīr}.) A pavement? (\textbf{\hi{फ़ुटपाथ}}, \emph{fuṭpāth}.) To take a stroll? (\textbf{\hi{टहलना}}, \emph{ṭahalnā}.) To pant? (\textbf{\hi{हाँफना}}, \emph{hā̃phnā}.) Traffic? (\textbf{\hi{ट्रैफ़िक}}, \emph{ṭraifik}.) A traffic jam? (\textbf{\hi{जाम}}, \emph{jām}.) A motorcycle? (\textbf{\hi{मोटरसाइकिल}}, \emph{moṭarsāikil}.) To crawl? (\textbf{\hi{रेंगना}}, \emph{reṅgnā}.) To drive (animals)? (\textbf{\hi{हाँकना}}, \emph{hā̃knā}.) Speed (formal)? (\textbf{\hi{गति}}, \emph{gati}.) A pothole? (\textbf{\hi{गड्ढा}}, \emph{gaḍḍhā}.) A conductor? (\textbf{\hi{कंडक्टर}}, \emph{kaṇḍakṭar}.) To bounce? (\textbf{\hi{उछलना}}, \emph{uchalnā}.) To be cut? (\textbf{\hi{कटना}}, \emph{kaṭnā}.)
\end{tcolorbox}
